% Part: counterfactuals
% Chapter: minimal-change-semantics
% Section: transitivity

\documentclass[../../../include/open-logic-section]{subfiles}

\begin{document}

\olfileid[tr]{cnt}{min}{tra}

\olsection{Geçişlilik}

Maddi koşul için zincir kuralı geçerlidir:
$!A\lif !B,!B\lif !C\Entails !A\lif !C$. Başka bir deyişle maddi koşul
geçişlidir. Aynısı karşıolgusallar için de doğru mudur? Stalnaker'ın şu
örneğini düşünün:
\begin{quote}
  J.~Edgar Hoover Rus olarak doğmuş olsaydı komünist olurdu.

  J.~Edgar Hoover komünist olsaydı hain olurdu.

  Öyleyse J.~Edgar Hoover Rus olarak doğmuş olsaydı hain olurdu.
\end{quote}
Hoover (gerçekte doğduğu zamanda) ABD'de değil de Rusya'da doğmuş olsaydı
Sovyetler Birliği'nde büyür ve komünist olurdu (böyle varsayalım). Dolayısıyla
birinci öncül doğrudur. Tek başına ele alındığında ikinci öncül de doğrudur.
Fakat sonuç yanlıştır: Hoover SSCB'de doğmuş olsaydı büyük olasılıkla ateşli
bir komünist olur, (ülkesine) ihanet etmezdi. Bu sezgisel doğruluk değeri
ataması Stalnaker--Lewis açıklamasınca da desteklenir. Bizim dünyamıza en
yakın, tek değişikliğin Hoover'ın doğum yeri olduğu olası dünya, Hoover'ın
SSCB'nin iyi bir yurttaşı olarak yetiştiği dünyadır. Bu, birinci öncülün ve
sonucun önbileşeninin doğru olduğu en yakın olası dünyadır; bu dünyada Hoover
Komünist Partinin sadık bir üyesidir ve dolayısıyla hain değildir. Buna
karşılık ikinci öncülü değerlendirmek için başka bir dünyaya bakmalıyız:
Hoover'ın komünist olduğu en yakın dünya, onun ABD'de doğup saf değiştirdiği
ve böylece hain olduğu dünyadır.\footnote{Elbette örneğin gücünü görmek için
  bazı metafizik ve siyasal varsayımları kabul etmemiz gerekir; örneğin
  Hoover'ın Rus ebeveynlerden doğmasının mümkün olduğu ya da 1950'lerde
  ABD'deki komünistlerin ülkelerine ihanet etmiş sayıldığı gibi.}

\begin{prob}
  Karşıolgusalların geçişsizliğine ikna edici, sezgisel bir örnek bulun.
\end{prob}

\begin{ex}\ollabel{ex:trans-counterex}
  Küre semantiği çıkarımı geçersiz kılar; yani
  $p\cif q,q\cif r\Entails/ p\cif r$. $W=\{w,w_1,w_2\}$,
  $O_w=\{\{w\},\{w,w_1\},\{w,w_1,w_2\}\}$,
  $V(p)=\{w_2\}$, $V(q)=\{w_1,w_2\}$ ve $V(r)=\{w_1\}$ olan
  $\mModel{M}=\tuple{W,O,V}$ modelini ele alın. $p$'yi kabul eden
  $S=\{w,w_1,w_2\}$ küresi vardır ve $p\lif q$ bu kürenin bütün
  dünyalarında doğrudur; dolayısıyla $\mSat{M}{p\cif q}[w]$. Ayrıca $q$'yu
  kabul eden $S'=\{w,w_1\}$ küresi vardır ve $\mSat{M}{q\lif r}$ bu
  kürenin bütün dünyalarında doğrudur; dolayısıyla $\mSat{M}{q\cif r}[w]$.
  Buna karşılık $p$'yi kabul eden $\{w,w_1,w_2\}$ küresi,
  $\mSat/{M}{p\lif r}[w_2]$ olan bir dünya, yani~$w_2$'yi içerir.
\end{ex}

\begin{prob}
\olref[cnt][min][tra]{ex:trans-counterex}'deki karşı örneğe karşılık gelen
küre şemasını çizin.
\end{prob}

\begin{prob}
  \olref[cnt][min][tra]{ex:trans-counterex}'de $w_2$, Hoover'ın Rusya'da
  doğduğu, komünist olduğu ve hain olmadığı dünyadır; $w_1$ ise Hoover'ın
  ABD'de doğduğu, komünist ve hain olduğu dünyadır. Bu modelde $w_1$, $w$'ye
  $w_2$'den daha yakındır. Bu zorunlu mudur? Hoover'ın Rusya'da doğmasının,
  komünist olmasından daha uzak bir olasılık olduğunu varsaymayan bir karşı
  örnek verebilir misiniz?
\end{prob}

\end{document}
